(1) 在 $\sigma_z=\operatorname{diag}(1,-1)$ 的计算基下，$\sigma_z\rho\sigma_z$ 保留两个对角元但改变两个非对角元的符号。因此 $\dot\rho_{00}=0,\dot\rho_{11}=0$，而 $\dot\rho_{01}=-\gamma\rho_{01}$、$\dot\rho_{10}=-\gamma\rho_{10}$，此动力学只消除相干性。
(2) 对矩阵元直接积分得 $\rho(t)=\begin{pmatrix}\rho_{00}(0)&e^{-\gamma t}\rho_{01}(0)\\e^{-\gamma t}\rho_{10}(0)&\rho_{11}(0)\end{pmatrix}$。总迹不变，初始对角元都保持常数；非对角元随指数衰减，但不会凭空改变各能级的占据概率。
(3) 初态 $|+x\rangle=(|0\rangle+|1\rangle)/\sqrt2$ 对应 $\rho(0)=\frac12\begin{pmatrix}1&1\\1&1\end{pmatrix}$。演化后 $\rho_{01}=\rho_{10}=e^{-\gamma t}/2$，故 $\langle\sigma_x\rangle=\operatorname{tr}(\rho\sigma_x)=e^{-\gamma t}$，$\langle\sigma_y\rangle=0$，横向退相干时间 $T_2=1/\gamma$。
(4) 对上述密度矩阵，$\operatorname{tr}(\rho^2)=\frac12(1+e^{-2\gamma t})$，从纯态值一单调降至 $1/2$。极限 $\rho(\infty)=I/2$ 是最大混合态；本模型没有能量交换型跃迁项，因此退相干并不表示人口占据向某个低能态弛豫。
